\section{Discussion}
\label{sec:discussion}

\subsection{A null confirmatory result that is not a null finding}

Ten pre-registered tests, none surviving correction --- and the reason is measurable rather than
mysterious. The design was sized on a standard deviation borrowed from the preceding
chapter~\citep{companion2026baseline}, and that borrowing is where the study earns its keep: the
quantity borrowed turns out not to be a property of the model, so no care at design time would have
produced a different number of binaries. Where the pre-registration recorded 84\% run-to-run noise and
a floor of thirteen binaries, the same nominal contrast here gives 0.5\% and 668. Not an underestimate
--- an inversion, and the lever the previous chapter said would work is the one that buys nothing here.

This is the finding we did not go looking for, and \S\ref{sec:rq6} places it against the literature on
borrowed effect sizes. The consequence for a designer belongs here: inheriting that calibration would
not merely have left them under-powered, it would have sent their budget to the axis that buys
nothing. The rival account for the decomposition --- that
the granularity of the metric produces the spread on its own --- is bounded in
\S\ref{sec:granularita} by the two conditions it would and would not have to explain, with the
figures for each.

\subsection{Why the census carries the weight}

The confirmatory arm and the census differ in what it takes to overturn them. Every number in
\S\ref{sec:results} is a point estimate with an interval that a larger study could move; nothing in
\S\ref{sec:census} is an estimate, since an endpoint either returned that message or it did not and the
artifact carries the request that produced it. A referee who disbelieves the census can re-run the probe;
one who disbelieves the contrasts has to re-run 5{,}809 trials.


\subsection{When to discard a batch, and when not to}

Two batches were invalidated and re-collected rather than patched, and the rule separating them from the
one we kept is reusable. \textbf{Discard when the contaminating channel is systematic and asymmetric
between the compared arms; report and bound it otherwise.}

The larger invalidation --- 2{,}390 measurements --- meets that test exactly: the algorithm's name reached
the model through three separate channels and the two arms encountered one at different rates, so left in
place it would have turned part of the transport effect into a measure of how often each transport stumbled
onto the answer. No downstream analysis separates a contaminant that correlates with the
manipulated variable.

The shared working directory (\S\ref{sec:threats}) fails the same test in the other direction:
non-directional noise, some thirty times in the original batch's 5{,}805 measurements, with no asymmetry
between arms.
Discarding it would have been a gesture rather than a correction, so we bounded it, probed it, and
re-collected on isolated workspaces. Two practices made the rule applicable rather than rhetorical:
channels were enumerated per channel and not per intention --- the third was found because we listed every
place the model could read --- and both batches remain in the artifact with their recomputation.

\subsection{What an evaluation should report}

What follows amends an existing checklist rather than proposing a new one: the community guidelines already
require the harness to be described once it goes beyond bare API calls~\citep{baltes2026guidelines}, and each
field below either sharpens that requirement or extends it to one they do not name. Each costs a line in a
method section, and Table~\ref{tab:checklist} states them as fields because a recommendation phrased as prose
is one a reader agrees with and does not apply. The first carries the argument of \S\ref{sec:runtime}: a
model name does not identify a measurable system, the same nominal model on two clouds being two apparatuses
that differ in score stability as well as in score.

\begin{table}[t]
\centering
\caption{What an agentic evaluation should declare. The first four sharpen requirements that
already exist~\citep{baltes2026guidelines}; the last two have no counterpart in current
checklists.}
\label{tab:checklist}
\footnotesize
\setlength{\tabcolsep}{4pt}
\begin{tabular}{@{}p{0.26\linewidth}p{0.66\linewidth}@{}}
\toprule
field & what it means \\
\midrule
transport & how a tool call is encoded: native function calling, a textual protocol, typed
  code stubs --- and whether more than one call per turn is admitted \\
endpoint & which service served the weights, not only the model name \\
model and version & the identifier the endpoint accepts, including any regional prefix \\
configuration & temperature, token limits, turn budget, seed where honoured \\
\textbf{models excluded} & which models were considered and not measured \\
\textbf{reason for each} & the endpoint's own message, verbatim, or the policy that applied \\
\bottomrule
\end{tabular}
\end{table}

\subsection{Five propositions that should outlive this task}
\label{sec:proposizioni}

The effect sizes here belong to one task family, four models and two clouds, and we claim nothing beyond
that. What we offer beyond it are five propositions about how experiments of this kind behave, each
instantiated by this study rather than proved in general.

\begin{enumerate}
  \item The transport belongs beside temperature in a method section.
  \item An endpoint can produce missing experimental units, whose cost \S\ref{sec:bound} bounds.
  \item Variance components can be apparatus-specific, so an inherited calibration must be recomputed on
        one's own data as soon as the first cell exists.
  \item A pre-flight must exercise realistic conversational states: two of our eight mechanisms fire only
        in a state a single-call probe never reaches.
  \item Metric granularity limits what a variance decomposition can say. We report the floor and ceiling
        shares beside ours, and did not size the tests per task.
\end{enumerate}

\noindent
The artifact carries each with its argument.


